%!TeX root=../main.tex

\chapter{Implementação do protótipo}

\label{cap:implementacao}

O Capítulo~\ref{cap:solucao} disse o que cada componente faz e por quê; este
diz como cada um foi construído, para quem quiser ler o código ou reproduzir
o protótipo. O código está em dois repositórios: o \textit{fork} do kw, no
ramo de telemetria \citep{kwfork_telemetry}, e o repositório deste trabalho,
com o kw\_tele, a configuração do Collector, o backend próprio e a
demonstração \citep{tccrepo}. As mudanças no kw ficam nesse ramo e não foram
propostas ao projeto. O capítulo segue o caminho do dado, do kw ao backend, e
termina nos testes.

\section{As mudanças no kw}
\label{sec:impl-kw}

A telemetria entra no kw por um único ponto: o script principal carrega o
disparador antes de escolher o comando a rodar (Programa~\ref{prog:entrada}).
O arquivo é lido com o \texttt{.} do Bash, e não com a função
\texttt{include} do kw, porque a função cria um processo a mais e esse trecho
roda em todo comando. O ramo acrescenta ao kw cerca de 3,6 mil linhas em 30 arquivos,
das quais 2,7 mil são testes; do resto, a maior parte é o disparador, com
pouco mais de 400 linhas, e o comando \texttt{kw telemetry}, com pouco mais de
100 \citep{kwfork_telemetry}.

\begin{program}[htbp]
  \centering
\begin{lstlisting}[language=bash, style=wider]
# Plain `.` instead of `include`, which forks `realpath`: this runs on every
# command. A partial installation just skips the gate.
if [[ -r "${KW_LIB_DIR}/lib/telemetry.sh" ]]; then
  . "${KW_LIB_DIR}/lib/telemetry.sh" --source-only &&
    telemetry_gate "$action"
fi
\end{lstlisting}
  \caption{O ponto de entrada do disparador no script principal do kw, antes
    da escolha do comando; a última chamada foi quebrada em duas linhas para
    caber na página \citep{kwfork_telemetry}.}
  \label{prog:entrada}
\end{program}

No banco, o arquivo do esquema, \texttt{database/kwdb.sql}, ganhou a coluna de
marca de envio na tabela de eventos, a tabela da lista de módulos e uma
\textit{view} de inserção com um gatilho que grava o evento, a execução e os
módulos numa instrução só. É o gatilho que torna a escrita atômica: uma
segunda chamada ao \texttt{sqlite3} não saberia o identificador do evento que
a primeira acabou de criar. A gravação dos módulos é chamada pelo
\texttt{build.sh} e pelo \texttt{deploy.sh}, com sucesso ou falha. A
verificação do esquema, para os bancos já instalados, acrescenta a coluna e um
índice dos eventos pendentes e reaplica o arquivo do esquema se a \textit{view} da lista de módulos faltar. Por fim, os comandos do kw passaram a esperar até 5 segundos quando o
banco está ocupado, para que a escrita do kw não falhe enquanto o kw\_tele
marca linhas.

O disparador segue a ordem da Figura~\ref{fig:disparador}. Ele lê a configuração e o estado com o próprio Bash e, no caminho comum, só
consulta o banco para contar os eventos pendentes; a verificação do esquema
roda apenas antes de perguntar o consentimento ou de disparar o kw\_tele. Para
disparar, ele usa \texttt{setsid}, que separa o kw\_tele do terminal, e
\texttt{timeout}, que limita a sua duração. Ele nunca
escreve na saída do comando nem muda o código de saída dele; a pergunta de
consentimento vai para a saída de erro.

O comando \texttt{kw telemetry} reúne o
que o usuário faz com a telemetria: ver o que foi gravado e enviado, enviar
ou ver o próximo lote sem enviar, aceitar ou recusar, e mudar a configuração;
o ramo acrescenta a página de manual correspondente. A configuração,
\texttt{etc/telemetry.config}, tem as chaves do comportamento descrito na
Seção~\ref{sec:cliente}: o intervalo entre envios, o tempo mínimo entre
tentativas, o tamanho do lote, o endereço e a lista de exclusão, que vem vazia. O
consentimento fica preso ao endereço aceito: se o endereço muda à mão, o kw
pergunta de novo.

O \texttt{setup.sh} do kw instala o kw\_tele. O caminho previsto é baixar a
versão publicada no repositório deste trabalho; enquanto não há versão
publicada, o kw\_tele é construído a partir de uma cópia local do repositório,
indicada por uma variável de ambiente. Numa reinstalação, o kw\_tele antigo é mantido até o novo ser instalado. Se a instalação falha, o \texttt{setup.sh}
avisa e continua, e o kw funciona sem telemetria, porque o disparador vê que o
kw\_tele não está instalado. A única dependência nova declarada para as
distribuições é o Python.
